\documentclass[../../../include/open-logic-section]{subfiles}

\begin{document}

\olfileid{sth}{spine}{rank}
\olsection{Rank}

Saiki tataran-tataran wis ditetepake minangka
$V_\alpha$, lan kita ngerti yen saben himpunan
iku himpunan bagean saka sawijining tataran.
Mula kita bisa netepake \emph{rank} sawijining
himpunan. Miturut intuisi, rank $A$ iku indeks
tataran paling awal kang unsur-unsure wis cukup
kanggo mbentuk $A$. Luwih persis:

\begin{defn}\ollabel{defnsetrank}
Kanggo saben himpunan $A$, $\setrank{A}$ iku
ordinal paling cilik $\alpha$ kang nyukupi
$A \subseteq V_\alpha$.
\end{defn}
\begin{prop}\ollabel{ranksexist}
	$\setrank{A}$ ana kanggo saben $A$.
\end{prop}
\begin{proof}
	Dipasrahake minangka latihan.
\end{proof}
\begin{prob}
	Buktikna \olref[sth][spine][rank]{ranksexist}.
\end{prob}
\noindent 
Amarga rank-rank terurut becik, kita bisa
mbuktekake sawetara asil kang wigati:
\begin{prop}\ollabel{valphalowerrank}
Kanggo saben ordinal $\alpha$,
$V_\alpha = \Setabs{x}{\setrank{x} \in \alpha}$.
\end{prop}

\begin{proof}
Yen $\setrank{x} \in \alpha$, mula
$x \subseteq V_{\setrank{x}} \in V_\alpha$,
saengga $x \in V_\alpha$ amarga $V_\alpha$
poten (kanthi nggunakake
\olref[Valphabasic]{Valphabasicprops} kaping pirang-pirang).
Kosokbaline, yen $x \in V_\alpha$, mula
$x \subseteq V_\alpha$, dadi
$\setrank{x} \leq \alpha$; induksi transfinit
prasaja banjur nuduhake yen
$\setrank{x} \in \alpha$. Ing kasus penerus,
kabeh unsur wis ana ing tataran sadurunge;
ing kasus limit, himpunane wis ana ing
sawijining tataran luwih awal.
\end{proof}
\begin{prob}
	Rampungna induksi transfinit prasaja kang
	kasebut ing \olref[sth][spine][rank]{valphalowerrank}.
\end{prob}

\begin{prop}\ollabel{rankmemberslower}
Yen $B \in A$, mula $\setrank{B} \in \setrank{A}$.
\end{prop}

\begin{proof}
$A \subseteq V_{\setrank{A}} = \Setabs{x}{\setrank{x} \in
\setrank{A}}$ miturut \olref{valphalowerrank}.
\end{proof}
\noindent
Kanthi kasunyatan iki, kita bisa mbuktekake
asil kang ngidini kita mbuktekake pratelan
babagan \emph{kabeh himpunan} liwat sawijining
wujud induksi:

\begin{thm}[Skema Induksi-$\in$]
Kanggo saben formula $\phi$:
\[
	\forall A((\forall x \in A)\phi(x) \lif \phi(A)) \lif \forall A \phi(A).
\]
\end{thm}

\begin{proof}
Kita bakal mbuktekake kontrapositif. Upamane
$\lnot \forall A \phi(A)$. Miturut Induksi
Transfinit (\olref[ordinals][basic]{ordinductionschema}),
ana sawijining kang dudu-$\phi$ kanthi rank
paling cilik; yaiku\ ana $A$ kang nyukupi
$\lnot \phi(A)$ lan
$\forall x(\setrank{x} \in \setrank{A}
\lif \phi(x))$. Saiki yen $x \in A$, mula
$\setrank{x} \in \setrank{A}$ miturut
\olref{rankmemberslower}, saengga $\phi(x)$;
yaiku\ $(\forall x \in A)\phi(x) \land
\lnot \phi(A)$.
\end{proof}\noindent Pratelan kang kuwat iki
bisa diterangake kanthi ora formal mangkene.
$\phi$ diarani \emph{herediter} yen lan mung yen:
yen kabeh !!{element} sawijining himpunan nyukupi
$\phi$, himpunan iku dhewe uga nyukupi
$\phi$. Mula Induksi-$\in$ nerangake yen
$\phi$ iku herediter, saben himpunan nyukupi
$\phi$.

Kanggo mungkasi rembugan rank (sauntara iki),
kita bakal mbuktekake sawetara pratelan kang
wis bola-bali diwenehi pratandha sadurunge.

\begin{prop}\ollabel{ranksupstrict}
$\setrank{A} = \supstrict_{x \in A}\setrank{x}$.
\end{prop}

\begin{proof}
Tetepna $\alpha = \supstrict_{x \in A}\setrank{x}$.
Miturut \olref{rankmemberslower},
$\alpha \leq \setrank{A}$. Nanging, yen
$x \in A$, mula $\setrank{x} \in \alpha$,
saengga $x \in V_\alpha$ miturut
\olref{valphalowerrank}, lan mula
$A \subseteq V_\alpha$, yaiku
$\setrank{A} \leq \alpha$. Dadi
$\setrank{A} = \alpha$.
\end{proof}

\begin{cor}\ollabel{ordsetrankalpha}
Kanggo saben ordinal $\alpha$,
$\setrank{\alpha} = \alpha$.
\end{cor}

\begin{proof}
Kanggo induksi transfinit, upamane
$\setrank{\beta} = \beta$ kanggo saben
$\beta \in \alpha$. Saiki
$\setrank{\alpha} = \supstrict_{\beta \in
\alpha}\setrank{\beta} = \supstrict_{\beta \in
\alpha}\beta = \alpha$ miturut
\olref{ranksupstrict}.
%	Elinga dhisik yen $\setrank{\alpha} \neq \beta$
%	kanggo saben $\beta \in
%	\alpha$. Yen ora,
%	kita bakal duwe $\beta = \setrank{\beta} \in
%	\setrank{\alpha} = \beta$ miturut
%	\olref{rankmemberslower}, yaiku
%	$\beta \in \beta$, sawijining kontradiksi.
%
%	Saiki elinga yen $\alpha \subseteq V_\alpha$.
%	Awit $\alpha =
%	\Setabs{\beta}{\beta \in \alpha}$ miturut
%	\olref[sfr][ordinals][basic]{ordissetofsmallerord}.
%	Lan yen $\beta \in \alpha$, mula
%	$\beta \subseteq V_{\setrank{\beta}} =
%	V_\beta \in V_\alpha$, saengga
%	$\beta \in V_\alpha$, miturut
%	\olref[sfr][spine][Valphabasic]{Valphabasicprops}
%	kaping pindho.
\end{proof}

Pungkasan, iki bukti cekak asil kang dijanjekake
ing pungkasan \olref[foundation]{sec}, yaiku
$\ZFminus$ mbuktekake implikasi
$\emph{Regularitas} \Rightarrow \emph{Landhesan}$.
(Elinga yen gagasan ``rank'' lan
\olref{rankmemberslower} bisa digunakake ing
bukti iki, amarga---kaya kang kasebut ing
wiwitan bagean iki---loro-lorone bisa
ditetepake kanthi $\ZFminus + \text{Regularitas}$.)

\begin{prop}[ing $\ZFminus + \text{Regularitas}$]\ollabel{zfminusregularityfoundation} Landhesan lumaku.
\end{prop}

\begin{proof}
Tetepna $A \neq \emptyset$, lan pilih
$B \in A$ kanthi rank paling cilik. Yen
$c \in B$, mula $\setrank{c} \in \setrank{B}$
miturut \olref{rankmemberslower}, saengga
$c \notin A$ amarga pilihan $B$.
\end{proof}

\end{document}
